% TOP_PROBLEM: {"id": 20000976, "title": "Periodic Laplacian certificates for robustness and explosion", "category_id": 2, "category": {"id": 2, "name": "combinatorics", "display_name": "Combinatorics", "description": "Counting problems, graph theory, discrete structures.", "slug": "combinatorics", "order_index": 2, "created_at": "2026-07-31T15:26:25.671Z"}, "status": "partially_solved", "problem_number": "AIM-COMBINATORICS-0101", "set_id": 14}
% FIRST_POSTED: 2026-10-01
% ENTRY_KIND: statement_only
% UnsolvedMath ID 20000976; source catalogue code AIM-COMBINATORICS-0101
% !TeX program = lualatex
% Definition and sources only; this file is not an LLM solution attempt.
\documentclass[11pt,a4paper]{article}
\usepackage[margin=25mm]{geometry}
\usepackage{fontspec}
\setmainfont{lmroman10-regular.otf}[BoldFont=lmroman10-bold.otf,
 ItalicFont=lmroman10-italic.otf,BoldItalicFont=lmroman10-bolditalic.otf]
\usepackage{amsmath,amssymb,mathtools}
\usepackage{unicode-math}
\setmathfont{latinmodern-math.otf}
\AtBeginDocument{\let\setminus\smallsetminus}
\usepackage{xurl,hyperref}
\hypersetup{colorlinks=true,urlcolor=blue}
\setcounter{secnumdepth}{0}
\setlength{\parindent}{0pt}
\setlength{\parskip}{5pt}
\setlength{\emergencystretch}{3em}
\begin{document}
\section*{Periodic Laplacian certificates for robustness and explosion}\label{20000976}
{\small UnsolvedMath ID 20000976 · AIM-COMBINATORICS-0101 · Combinatorics · AIM Workshop Problem Lists}
\subsection{Definitions and mathematical statement}
On the square lattice $\mathbb Z^2$, a sandpile configuration assigns a nonnegative integer number of grains to every vertex.
A vertex with at least four grains may topple legally: it loses four grains and sends one to each of its four nearest neighbors.
A configuration is stable when every height belongs to $\{0,1,2,3\}$.
Call a configuration stabilizable if there exists a finite or countably infinite legal toppling sequence in which every vertex topples only finitely many times and whose limiting configuration is stable.
This allows activity at infinitely many different vertices; it does not require a globally finite total number of topplings.

The input is a stable periodic configuration $c$ specified by two linearly independent integer period vectors and its finitely many heights on representatives for their lattice quotient.
Let $\delta_0$ be one grain at the origin and zero elsewhere.
Is there an algorithm, halting on every such finite input, that decides the assertion
\[
 \exists n\in\mathbb Z_{\geq1}\quad
           c+n\delta_0\text{ is not stabilizable}?
\]
The universal alternative is that $c+n\delta_0$ stabilizes for every finite $n$.
The period lattice and its size are unrestricted parts of the input, rather than fixed in advance.
The question concerns decidability of the existence of some destabilizing addition; direct simulation of a chosen addition alone does not provide a halting decision procedure for both alternatives.
No finite boundary or absorbing sink is present in this infinite-lattice model.
\subsection{Short English statement}
Can one decide from a periodic square-lattice background whether adding finitely many grains at the origin can ever cause nonstabilization?
\subsection{Sources}
\begin{itemize}
\item[S1] ULAM.AI and the UnsolvedMath Contributors, supplied problems.json, record 20000976 (AIM-COMBINATORICS-0101), AIM Workshop Problem Lists. Catalogue identity and source transcription; CC BY 4.0.
\url{https://huggingface.co/datasets/ulamai/UnsolvedMath}.
\item[S2] American Institute of Mathematics, Generalizations of chip-firing and the critical group, item 11. Original workshop question underlying AIM-COMBINATORICS-0101.
\url{https://aimath.org/pastworkshops/chipfiringproblems.pdf}.
\end{itemize}
\end{document}
